\section{Information-theoretic bounds and experimental design}\label{sec:info}

\subsection{A sandwich for $\alpha$: the Chernoff upper bound and the union lower bound}

\begin{proposition}[Chernoff upper bound, with an explicit finite-$d$ form]\label{prop:genie}
Let $d\ge3$ be odd, consider the rotated surface-code family (conventions of \texttt{lcd.codes.RotatedSurfaceCode}) and noise $\Nn_{\rm cc}(p,e)$ with $0\le p\le\tfrac34$, $0\le e\le1$, $(p,e)\ne(0,0)$. Then
\[
\alpha_+(p,e)\;\le\;\ln\frac1{B(p,e)},\qquad B(p,e)=e+(1-e)\Big(2\sqrt{\tfrac{p(1-p)}3}+\tfrac{2p}3\Big).
\]
More precisely, for every $d$ we have $\err_d\ge\varepsilon_{\rm genie}(d)$, where
$\varepsilon_{\rm genie}(d)=\tfrac12\bigl(1-\mathrm{TV}(P_1^{\otimes d},Q_1^{\otimes d})\bigr)$,
$P_1$ is the joint law of a single qubit's $(a,f)$ and $Q_1=P_1(\cdot+X)$ is the shifted law, and
\[
\varepsilon_{\rm genie}(d)\;\ge\;\tfrac12\,B^d\exp\!\Big(-\tfrac12\sqrt{d\,\sigma^2}\Big),\qquad
\sigma^2=\frac{2(1-e)\sqrt{(1-p)p/3}}{B}\,\ln^2\!\frac{3(1-p)}{p}\quad(\sigma^2=0\text{ if }p=0).
\]
\hfill\st{PROVED}
\end{proposition}
\begin{proof}
Take the $d$ qubits $T$ of an interior column $c_0$, $1\le c_0\le d-2$, and let $\bar X_T$ be the $X$-type operator on this column.

\emph{Step 1: $\ker\sigma\cap\mathcal P_T=\{I,\bar X_T\}$, and $\bar X_T\notin S$.} A plaquette with top-left corner $(r,c)$ covers $(r,c),(r,c+1),(r+1,c),(r+1,c+1)$ and is of $Z$ type iff $r+c$ is odd.
For each $r\in\{0,\dots,d-2\}$ the two bulk plaquettes with corners $(r,c_0-1)$ and $(r,c_0)$ exist (because $1\le c_0\le d-2$) and have opposite types; each meets column $c_0$ exactly in $\{(r,c_0),(r+1,c_0)\}$.
Let $(x,z)$ be supported on $T$ with zero syndrome. The $Z$-type plaquette of the pair forces $x_r=x_{r+1}$ for every $r$, so $x\in\{0,\mathbf 1_T\}$; the $X$-type one forces $z_r=z_{r+1}$, so $z$ is constant.
The weight-$2$ $X$-type checks on the top edge have corners $(-1,c)$ with $c$ odd and cover $(0,c),(0,c+1)$; exactly one of them ($c=c_0$ if $c_0$ is odd, $c=c_0-1$ if $c_0$ is even) contains $(0,c_0)$, and it meets $T$ only there. Hence $z_0=0$ and $z=0$.
Conversely $\bar X_T$ has zero syndrome: bulk $Z$ plaquettes meet the column in $0$ or $2$ qubits and the $Z$-type boundary checks lie on columns $0$ and $d-1$. It anticommutes with the $Z$-type logical on row $0$ (overlap $1$), so $\bar X_T\notin S$ and it flips the label.
(Check K1 of \texttt{checks/completion\_checks.py} verifies the statement over $\F$ for every interior column, $d=3,\dots,41$.)

\emph{Step 2: the genie.} The genie knows $F$, $E_{T^c}$, and $\sigma(E)$, which is equivalent to knowing $\sigma(E_T)$. By Step 1, $E_T$ can only be $a$ or $a+\bar X_T$,
and the two have different labels. The genie's Bayes error is $\tfrac12\sum_{a,f_T}\min\big(P_T(a,f_T),P_T(a+\bar X_T,f_T)\big)=\varepsilon_{\rm genie}$.
The real observation $(\sigma(E),F)$ is a coarse-graining (F1) of the genie's observation, so $\err_d\ge\varepsilon_{\rm genie}$ (Theorem~\ref{thm:monotone}).

\emph{Step 3: the explicit bound.} Write $P=P_1^{\otimes d}$, $Q=Q_1^{\otimes d}$ and $\Lambda=\ln(P/Q)$ where $PQ>0$. Since $\min(P,Q)=\sqrt{PQ}\,e^{-|\Lambda|/2}$ (both sides vanish where $PQ=0$) and $\sum_x\sqrt{P(x)Q(x)}=B^d$
($B$ is the Bhattacharyya coefficient of $P_1$ and $Q_1$, \S\ref{sec:exchange}),
\[
\varepsilon_{\rm genie}=\tfrac12B^d\,\E_R\big[e^{-|\Lambda|/2}\big],\qquad R:=\sqrt{PQ}/B^d=R_1^{\otimes d},\quad R_1=\sqrt{P_1Q_1}/B .
\]
Under $R$, $\Lambda=\sum_i\lambda(x_i)$ with $x_i$ i.i.d.\ $\sim R_1$. The involution $x\mapsto x+X$ exchanges $P_1$ and $Q_1$, so it preserves $R_1$ and maps $\lambda$ to $-\lambda$; hence $\E_{R_1}\lambda=0$ and $\E_R\Lambda^2=d\,\E_{R_1}\lambda^2=d\sigma^2$.
On an unflagged qubit $\lambda=\pm\ln\frac{3(1-p)}p$ for $a\in\{I,X\}$ (each of $R_1$-mass $(1-e)\sqrt{(1-p)p/3}/B$) and $\lambda=0$ for $a\in\{Y,Z\}$; on a flagged qubit $\lambda=0$; this gives the stated $\sigma^2$.
By Jensen's inequality (convexity of $\exp$) and Cauchy--Schwarz, $\E_R e^{-|\Lambda|/2}\ge e^{-\E_R|\Lambda|/2}\ge e^{-\sqrt{\E_R\Lambda^2}/2}=e^{-\sqrt{d\sigma^2}/2}$.
Finally $-\tfrac1d\ln\err_d\le\ln\tfrac1B+\tfrac{\ln2}d+\tfrac{\sigma}{2\sqrt d}$, whose $\limsup$ is $\ln(1/B)$.
Note that the flags live in the joint alphabet $(a,f)$ and the genie sees them, so at positions with $F=1$ the two hypotheses are \emph{identical},
contributing $e\cdot1$ to $B$; this is the origin of the $e$ term in $B$.
\end{proof}

\noindent The finite-$d$ bound replaces the earlier appeal to the asymptotic Chernoff theorem \cite[\S11.9]{CT}; its exponent is the Chernoff information $-\ln B$ (the minimum over $s$ of $b_s$ is at $s=\tfrac12$, \S\ref{sec:exchange}).
Check K2 compares the bound with the exact $\varepsilon_{\rm genie}(d)$ for $d\le61$ (ratio $\ge1.67$ on the tested grid) and the slope of $-\ln\varepsilon_{\rm genie}$ with $\ln(1/B)$ (within $2\%$ between $d=101$ and $201$).
For $d=3$, $\err\ge\varepsilon_{\rm genie}$ has been verified by exact computation (check C6); \texttt{checks/genie\_large\_d.py} tests it against the ML data of P1 for $d=5$--$13$.

\begin{proposition}[Union lower bound]\label{prop:lower}
Let $\Phi(x)=\limsup_{d\to\infty}\tfrac1d\ln\widetilde W_d(x)$ (Corollary~\ref{cor:product}). Then
$\alpha_-(p,e)\ge-\Phi\big(B(p,e)\big)$.\hfill\st{PROVED}
\end{proposition}
\begin{proof}
Corollary~\ref{cor:product}: $\err_d\le\tfrac12\widetilde W_d(B)$; take $-\tfrac1d\ln$ and then $\liminf$.
\end{proof}

\begin{corollary}[Sandwiched between two functions of $B$ only]\label{cor:sandwich}
$-\Phi(B(p,e))\le\alpha_-\le\alpha_+\le-\ln B(p,e)$.\hfill\st{PROVED}
\end{corollary}

The sandwich does not constrain the derivatives of $\alpha$, so it does not directly give upper and lower bounds on the exchange rate; its role is to give the testable hypothesis $H_B$
(Conjecture~\ref{conj:HB}): if $\alpha$ is a function of $B$ in a neighborhood, then $R_\alpha=R_B$.

\subsection{Reference values for the exchange rate}

The table below lists three quantities at the README work points, all computable from closed-form expressions. Here $c$ is the rigorous upper bound (Theorem~\ref{thm:rate});
$R_B$ and $R_{\rm hash}$ are \emph{heuristic references}, \textbf{not} theorems about the $R_\alpha$ of the surface code, and serve only to give the P1 numerics an order-of-magnitude expectation:
\[
R_{\rm hash}(p,e)=\frac{2-H_4(p)}{(1-e)\,\log_2\!\frac{3(1-p)}{p}},\quad
H_4(p)=-(1-p)\log_2(1-p)-p\log_2\tfrac p3,
\]
the slope of the level curves of the quantum hashing capacity $Q_{\rm hash}=1-2e-(1-e)H_4(p)$ (whose zero level curve gives $p\approx18.9\%$ at $e=0$
and $e=50\%$ at $p=0$).

\begin{center}
\begin{tabular}{cc ccc}
\toprule
$p_0$ & $e_0$ & $c=\frac{3/4-p_0}{1-e_0}$ (upper bound) & $R_B$ & $R_{\rm hash}$ \\
\midrule
0.02 & 0    & 0.7300 & 0.1784 & 0.2537 \\
0.02 & 0.02 & 0.7449 & 0.1820 & 0.2589 \\
0.02 & 0.05 & 0.7684 & 0.1877 & 0.2671 \\
0.04 & 0    & 0.7100 & 0.2212 & 0.2746 \\
0.04 & 0.02 & 0.7245 & 0.2257 & 0.2802 \\
0.04 & 0.05 & 0.7474 & 0.2328 & 0.2891 \\
0.06 & 0    & 0.6900 & 0.2444 & 0.2840 \\
0.06 & 0.02 & 0.7041 & 0.2494 & 0.2898 \\
0.06 & 0.05 & 0.7263 & 0.2573 & 0.2989 \\
\bottomrule
\end{tabular}
\end{center}
(These values are printed by check C7 of \texttt{theory/checks/exact\_small\_codes.py}. The same script also gives the finite-size exchange rates $R^{(3)}$ of the $d=3$ codes,
for example $0.54$ for the surface code ($p_0=0.02,e_0=0$); it satisfies $R^{(d)}\le c$ but lies above $R_B$ and $R_{\rm hash}$;
how the finite-size $R^{(d)}$ evolves with $d$ is exactly what P1 is to measure.)

\begin{remark}[Threshold secant]
Using the known thresholds $p_c\approx0.189$ (at $e=0$) and $e_c=0.5$ (at $p=0$), the secant slope of the zero level curve is $0.189/0.5\approx0.38$,
far below $c\approx0.75$. This is only a heuristic comparison of a secant (at threshold $\alpha=0$, and the local derivative at a work point is different), not a theorem.
As a rigorous consistency check: Corollary~\ref{cor:worth} gives $p_{\rm eq}\le3e/4$, so the pure-erasure threshold must satisfy
$e_c\ge\tfrac43p_c\approx0.25$, which is compatible with $e_c=0.5$.
\end{remark}

\subsection{Experimental design and detection limits}\label{sec:design}

This part corresponds to P2. Let $T$ be the total observation time (number of rounds) and let events be a Poisson process of rate $r$.

\begin{proposition}[Poisson lower bound]\label{prop:poisson}
Let $H_0$: no burst events; $H_1$: burst events form a Poisson process of intensity $r$ (of arbitrary shape), with all other layers independent and identical.
For any test $\varphi$ based on all data over the total time $T$, if the type-I error under $H_0$ is $\le a$, then the power under $H_1$ satisfies
\[
\Prob_1[\varphi=1]\;\le\;1-e^{-rT}(1-a).
\]
Therefore reaching power $1-\delta$ requires $T\ge\ln\!\big((1-a)/\delta\big)/r$. No method that exploits space-time patterns can beat the $1/r$ scaling.\hfill\st{PROVED}
\end{proposition}
\begin{proof}
Under $H_1$, with probability $e^{-rT}$ no event occurs at all, and then the data have the same law as under $H_0$. Hence
$\Prob_1[\varphi=1]\le(1-e^{-rT})+e^{-rT}\Prob_0[\varphi=1]\le1-e^{-rT}(1-a)$.
\end{proof}

\begin{proposition}[Chain of exponents for class discrimination]\label{prop:exponents}
Let events carry a label $a\in\mathcal A$ (the class), and let the class intensities under the two hypotheses be $r_a^{(0)},r_a^{(1)}$. With events fully observed,
the Stein exponent and the Chernoff exponent per unit time are
\[
D=\sum_a\Big[r^{(1)}_a\ln\frac{r^{(1)}_a}{r^{(0)}_a}-r^{(1)}_a+r^{(0)}_a\Big],\qquad
C=\max_{s\in[0,1]}\sum_a\Big[s\,r^{(1)}_a+(1-s)\,r^{(0)}_a-(r^{(1)}_a)^s(r^{(0)}_a)^{1-s}\Big]
\]
(the relative entropy and Chernoff information of Poisson processes). If the actual data $Z$ are the output of the event process through a detector channel $W$, then
\[
D_{\rm count}\ \le\ D_{\rm pattern}\ \le\ D_{\rm event}=D,\qquad C_{\rm count}\le C_{\rm pattern}\le C,
\]
where ``count'' means using only the total number of detected events and ``pattern'' means the optimal test using the full space-time pattern. Hence the ratio of the numbers of rounds needed to reach the same error satisfies
$T_{\rm count}/T_{\rm pattern}\to D_{\rm pattern}/D_{\rm count}\ (\delta\to0)$, bounded above by
$D_{\rm event}/D_{\rm count}$.\hfill The chain of inequalities is \st{PROVED} (data-processing inequality for relative entropy/R\'enyi divergences and the Poisson closed forms); the asymptotic form of the ratio of rounds is \st{PROVED} under the marking model (M) of \S\ref{sec:design-theorems} (Theorem~\ref{thm:stein-rounds}; the error held at a fixed level is the one under $H_1$)
\end{proposition}

\begin{problem}[O4: experimental-design theorem]\label{prob:O4}
Under a budget constraint choose the code distance $d$, the number of rounds, and the (adaptive) strategy so as to maximize $D_{\rm pattern}(d,\text{strategy})$ or the Chernoff exponent;
and characterize the dependence of $D_{\rm pattern}/D_{\rm count}$ on the detector channel $W$.\hfill\st{PLAN}
\end{problem}
\noindent The numerics of P2 estimate $D_{\rm count}$ and $D_{\rm pattern}$ separately (by Monte Carlo or analytically), report $T_{\rm count}/T_{\rm pattern}$, and compare with Proposition~\ref{prop:poisson}
(the lower bound $\ln(1/\delta)/r$ on the number of rounds for any method).

\subsection{Experimental design under the marking model (O4)}\label{sec:design-theorems}

This subsection proves the asymptotic part of Proposition~\ref{prop:exponents}, gives an exact formula for the gain of pattern methods over counting, and shows that adaptive choice among experiments does not improve the Stein exponent. Time is measured in rounds.

\paragraph{Assumption (M): marked Poisson observation.} Under hypothesis $h\in\{0,1\}$, (i) events of class $a$ occur as independent Poisson processes of rate $r^{(h)}_a$ per round; (ii) each event independently produces one observed \emph{mark} $z$ in a finite set $\mathcal Z$ with probability $W(z\mid a)$ ($\sum_zW(z\mid a)\le1$; the deficit is the probability that the event is not seen); (iii) spurious marks from the background form an independent Poisson process of intensity $b(z)$ per round, the same under both hypotheses; (iv) the marks of different events are observed separately and are assigned to the round in which the event occurs. Put
\[
\lambda_h(z)=b(z)+\sum_ar^{(h)}_aW(z\mid a),\qquad\Lambda_h=\sum_z\lambda_h(z),\qquad\pi_h=\lambda_h/\Lambda_h,
\]
and assume $\lambda_0(z)>0\Leftrightarrow\lambda_1(z)>0$ and $\Lambda_0,\Lambda_1>0$.
(M) idealizes the clustering of P2: (iv) fails when events overlap in space-time, which is rare when $r$ times the event volume is small.

\begin{lemma}[Divergences of the observed data]\label{lem:poisson-kl}
Under (M) the data of $T$ rounds consist of independent counts $N_z\sim{\rm Poisson}(T\lambda_h(z))$, and given the counts the rounds of the marks are i.i.d.\ uniform (independently of $h$). Hence the relative entropy of the data of $T$ rounds under $H_1$ relative to $H_0$ is $T D_{\rm pattern}$, and that of the total count $N=\sum_zN_z$ is $TD_{\rm count}$, where
\[
D_{\rm pattern}=\sum_z\Big[\lambda_1\ln\frac{\lambda_1}{\lambda_0}-\lambda_1+\lambda_0\Big](z),\qquad D_{\rm count}=\Lambda_1\ln\frac{\Lambda_1}{\Lambda_0}-\Lambda_1+\Lambda_0 .
\]
The same holds for the Chernoff information, with $C_{\rm pattern}=\max_s\sum_z[s\lambda_1+(1-s)\lambda_0-\lambda_1^s\lambda_0^{1-s}](z)$ and its count analogue. $D_{\rm event}$ and $C$ of Proposition~\ref{prop:exponents} are the case $W={\rm id}$, $b=0$.\hfill\st{PROVED}
\end{lemma}
\begin{proof}
By the marking and superposition theorems for Poisson processes, the marks of type $z$ form independent Poisson processes of rates $\lambda_h(z)$. The law of the data is the product over $z$ of the count laws times the conditional law of the rounds, which does not depend on $h$; relative entropy and Chernoff coefficients of products factorize, and $D({\rm Poisson}(\mu_1)\|{\rm Poisson}(\mu_0))=\mu_1\ln(\mu_1/\mu_0)-\mu_1+\mu_0$, $\sum_n{\rm Pois}_{\mu_1}(n)^s{\rm Pois}_{\mu_0}(n)^{1-s}=\exp(\mu_1^s\mu_0^{1-s}-s\mu_1-(1-s)\mu_0)$. $N$ is Poisson with mean $T\Lambda_h$.
\end{proof}

\begin{theorem}[What pattern information is worth]\label{thm:pattern-gain}
Under (M),
\[
D_{\rm pattern}=D_{\rm count}+\Lambda_1\,D(\pi_1\|\pi_0).
\]
Hence $D_{\rm pattern}\ge D_{\rm count}$, with equality iff $\pi_1=\pi_0$ (the marks carry no information: $\lambda_1(z)/\lambda_0(z)$ does not depend on $z$), and
\[
\frac{D_{\rm pattern}}{D_{\rm count}}=1+\frac{\Lambda_1D(\pi_1\|\pi_0)}{D_{\rm count}}.
\]
$D_{\rm pattern}$ does not increase when the marks are coarse-grained, $z\mapsto z'$ with a kernel $K(z'\mid z)$ (equivalently $W\mapsto KW$, $b\mapsto Kb$).
In the weak-signal limit $H_1$: one extra class $a^\ast$ of rate $\epsilon$ on top of $H_0$ (so $\lambda_1=\lambda_0+\epsilon w$, $w=W(\cdot\mid a^\ast)$, $\|w\|_1>0$),
\[
\lim_{\epsilon\to0}\frac{D_{\rm pattern}}{D_{\rm count}}=1+\chi^2\big(\hat w\,\big\|\,\pi_0\big),\qquad\hat w=w/\|w\|_1,
\]
so the gain of pattern methods for weak signals is exactly one plus the $\chi^2$-distance between the mark distribution of the sought events and that of all other marks.\hfill\st{PROVED}
\end{theorem}
\begin{proof}
Write $\lambda_h=\Lambda_h\pi_h$: $\sum_z\lambda_1\ln(\lambda_1/\lambda_0)=\Lambda_1\ln(\Lambda_1/\Lambda_0)+\Lambda_1\sum_z\pi_1\ln(\pi_1/\pi_0)$, and $\sum_z(\lambda_0-\lambda_1)=\Lambda_0-\Lambda_1$. Equality: $D(\pi_1\|\pi_0)=0$ iff $\pi_1=\pi_0$.
Coarse-graining: $D_{\rm pattern}=\sum_z\lambda_0\,f(\lambda_1/\lambda_0)$ with $f(t)=t\ln t-t+1$ convex, so the log-sum (Jensen) inequality gives monotonicity under any kernel, exactly as for $f$-divergences.
Weak signal: $D({\rm Pois}(\mu+\epsilon w)\|{\rm Pois}(\mu))=\tfrac{\epsilon^2w^2}{2\mu}+O(\epsilon^3)$ term by term, so $D_{\rm pattern}=\tfrac{\epsilon^2}2\sum_zw(z)^2/\lambda_0(z)+O(\epsilon^3)$ and $D_{\rm count}=\tfrac{\epsilon^2}2\|w\|_1^2/\Lambda_0+O(\epsilon^3)$; the ratio tends to $\Lambda_0\sum_z\hat w^2/\lambda_0=\sum_z\hat w^2/\pi_0=1+\chi^2(\hat w\|\pi_0)$.
\end{proof}

\begin{theorem}[Stein form of the ratio of rounds]\label{thm:stein-rounds}
Assume (M) and fix $a\in(0,1)$. Among tests based on $T$ rounds whose error under $H_1$ is at most $a$, let $\beta_T$ be the smallest error under $H_0$, and let $T^\ast(\delta)=\min\{T\in\mathbb N:\beta_T\le\delta\}$. If $0<D<\infty$, then
\[
\lim_{\delta\to0}\frac{T^\ast(\delta)}{\ln(1/\delta)}=\frac1D,
\]
with $D=D_{\rm pattern}$ for tests that use all data and $D=D_{\rm count}$ for tests that use the total count only. Hence $T^\ast_{\rm count}(\delta)/T^\ast_{\rm pattern}(\delta)\to D_{\rm pattern}/D_{\rm count}$, as stated in Proposition~\ref{prop:exponents}.
For the Bayes error with fixed priors the same holds with the Chernoff informations $C_{\rm pattern},C_{\rm count}$; holding the error under $H_0$ fixed instead gives the reversed divergences $D(P_0\|P_1)$, for which Theorem~\ref{thm:pattern-gain} holds with the roles of $0$ and $1$ exchanged.\hfill\st{PROVED}
\end{theorem}
\begin{proof}
By Lemma~\ref{lem:poisson-kl} the data of $T$ rounds are $T$ i.i.d.\ blocks with per-block divergence $D$. Stein's lemma \cite[Thm.~11.8.3]{CT} (for every fixed $a\in(0,1)$, with the roles of the hypotheses as stated) gives $-\tfrac1T\ln\beta_T\to D$.
$\beta_T>0$ for every $T$, since the two laws are mutually absolutely continuous, and $\beta_T$ is nonincreasing (a test may ignore rounds); hence $T^\ast(\delta)\to\infty$ as $\delta\to0$.
Given $\eta>0$ choose $T_\eta$ with $e^{-(D+\eta)T}\le\beta_T\le e^{-(D-\eta)T}$ for $T\ge T_\eta$. Then $T^\ast(\delta)\le\max\{T_\eta,\lceil\ln(1/\delta)/(D-\eta)\rceil\}$, and for $\delta$ so small that $T^\ast(\delta)\ge T_\eta$, $\delta\ge\beta_{T^\ast}\ge e^{-(D+\eta)T^\ast}$, i.e.\ $T^\ast\ge\ln(1/\delta)/(D+\eta)$. Let $\eta\to0$.
The count version is the same argument applied to the Poisson count process. The Bayes version uses the Chernoff theorem \cite[Thm.~11.9.1]{CT} in place of Stein's lemma.
\end{proof}

\begin{theorem}[Adaptive designs do not beat the best single experiment (Stein)]\label{thm:adaptive}
Let $\mathcal I$ be a finite set of experiments (for example: run a patch of distance $d$ for one round with a given readout and clustering scheme). Experiment $i$ costs $c_i>0$ and yields data with law $P_{h,i}$ under $H_h$, with $D_i:=D(P_{1,i}\|P_{0,i})<\infty$. An adaptive design chooses the next experiment as a function of all past data (and private randomness), stops when the total cost would exceed a budget $\mathcal B$, and then decides. Let $D^\ast=\max_iD_i/c_i$.
If the design has error at most $a$ under $H_1$, its error $\alpha_0$ under $H_0$ satisfies
\[
\ln\frac1{\alpha_0}\;\le\;\frac{\mathcal B\,D^\ast+\ln2}{1-a}.
\]
Repeating the single experiment $i^\ast=\arg\max_iD_i/c_i$ achieves $\liminf_{\mathcal B\to\infty}\tfrac1{\mathcal B}\ln\tfrac1{\alpha_0}\ge D^\ast$ at every fixed $a\in(0,1)$. Hence the optimal exponent per unit cost, in the limit $a\to0$, is $D^\ast$, and it is attained by a non-adaptive design.\hfill\st{PROVED}
\end{theorem}
\begin{proof}
Let $Q_h$ be the law of the whole transcript (choices and data) under $H_h$. The number of steps is at most $\mathcal B/\min_ic_i$. By the chain rule for relative entropy, and because the choice at step $t$ has the same conditional law given the past under both hypotheses,
$D(Q_1\|Q_0)=\E_{Q_1}\sum_tD_{i_t}\le D^\ast\,\E_{Q_1}\sum_tc_{i_t}\le D^\ast\mathcal B$.
Let $x=Q_1[\text{decide }1]\ge1-a$ and $y=Q_0[\text{decide }1]=\alpha_0$. Data processing to the binary decision gives $D(Q_1\|Q_0)\ge x\ln\tfrac xy+(1-x)\ln\tfrac{1-x}{1-y}\ge x\ln\tfrac1y-\ln2$, which is the bound. Achievability is Stein's lemma for $\lfloor\mathcal B/c_{i^\ast}\rfloor$ i.i.d.\ uses of $i^\ast$.
\end{proof}

\noindent\emph{Reading for P2.} Theorems~\ref{thm:pattern-gain}--\ref{thm:adaptive} settle the two-hypothesis part of O4 (Problem~\ref{prob:O4}) under (M): the design question reduces to computing $D_{\rm pattern}(i)/c_i$ for each candidate $i$ (distance, readout, clustering) and taking the largest; the gain over counting is $\Lambda_1D(\pi_1\|\pi_0)$, which for weak signals is governed by $\chi^2(\hat w\|\pi_0)$; and the Poisson bound (Proposition~\ref{prop:poisson}) remains the floor for every design.
What is \emph{not} settled: (a) composite hypotheses (unknown shape, unknown rates), where adaptive designs can help and the relevant exponent is a minimax over the composite alternative; (b) the strong converse at fixed $a$ for adaptive designs (Theorem~\ref{thm:adaptive} uses the weak converse; the strong converse for classical adaptive discrimination is known in the literature but not reproduced here); (c) the Chernoff (Bayes) exponent for adaptive designs; (d) the computation of $D_{\rm pattern}(d)$ for the surface code, which needs the mark channel $W$ of P2; (e) the regime where (iv) fails (overlapping events).
